\subsection{Example: linear representations in spaces of continuous functions}

Let G be a discrete group operating on the left on a set X. A complex function \(\chi\) on \(G \times X\) is called a multiplier if

\begin{enumerate}
\def\labelenumi{(\arabic{enumi})}
\tightlist
\item
\end{enumerate}

\[
\chi (e, x) = 1 \quad \text { for   all } x \in X;\tag{2}
\]

\[
\chi (s t, x) = \chi (s, t x) \chi (t, x) \quad \text { for   all } s, t \text { in } G, x \in X.
\]

It follows that

\[
\chi (t ^ {- 1}, t x) \chi (t, x) = 1 \quad \text { for   all } t \in \mathrm{G}, x \in \mathrm{X},\tag{3}
\]

and in particular \(\chi(t, x) \neq 0\) for all \(t \in \mathbf{G}\) , \(x \in \mathbf{X}\) .

For every complex function \(f\) defined on \(\mathbf{X}\) and every \(s \in \mathbf{G}\) , let \(\pmb{\gamma}_{\chi}(s)f\) be the complex function on \(\mathbf{X}\) defined by

\[
\left(\boldsymbol {\gamma} _ {\chi} (s) f\right) (x) = \chi (s ^ {- 1}, x) f (s ^ {- 1} x).\tag{4}
\]

Then \(\gamma_{\chi}(e)f = f\) and

\[
\begin{array}{r l} & {\big (\pmb {\gamma} _ {\chi} (s) \pmb {\gamma} _ {\chi} (s ^ {\prime}) f \big) (x) = \chi (s ^ {- 1}, x) \big (\pmb {\gamma} _ {\chi} (s ^ {\prime}) f \big) (s ^ {- 1} x)} \\ & {\qquad = \chi (s ^ {- 1}, x) \chi (s ^ {- 1}, s ^ {- 1} x) f (s ^ {- 1} s ^ {- 1} x)} \\ & {\qquad = \chi \big ((s s ^ {\prime}) ^ {- 1}, x) \big) f \big ((s s ^ {\prime}) ^ {- 1} x \big) = \big (\pmb {\gamma} _ {\chi} (s s ^ {\prime}) f \big) (x),} \end{array}
\]

thus \(\pmb{\gamma}_{\chi}\) is a linear representation of G. For \(\chi = 1\) , one recovers the endomorphisms \(\pmb{\gamma}(s)\) (Ch. VII, §1, No.~1, formula (3)).

Suppose now that G and X are locally compact, G operating continuously on X, and \(\chi\) continuous on \(G \times X\) . Then \(\mathcal{C}(X)\) and \(\mathcal{K}(X)\) are stable for the \(\boldsymbol{\gamma}_{\chi}(s)\) , whence linear representations of G in \(\mathcal{C}(X)\) and \(\mathcal{K}(X)\) which we shall again denote \(\gamma_{\chi}\) .

PROPOSITION 4. --- The linear representations \(\pmb{\gamma}_{\chi}\) of G in \(\mathcal{C}(\mathbf{X})\) and \(\mathcal{K}(\mathbf{X})\) are continuous.

The mapping \((s, f) \mapsto (s, \gamma(s)f)\) of \(\mathbf{G} \times \mathcal{C}(\mathbf{X})\) into \(\mathbf{G} \times \mathcal{C}(\mathbf{X})\) is continuous (No.~2, Lemma 3). On the other hand, the mapping \((s, f) \mapsto \chi(s, \cdot)f\) of \(\mathbf{G} \times \mathcal{C}(\mathbf{X})\) into \(\mathcal{C}(\mathbf{X})\) is continuous; for, if \(s\) tends to \(s_0\) in \(\mathbf{G}\) , then \(\chi(s, \cdot)\) tends to \(\chi(s_0, \cdot)\) uniformly on every compact subset of \(\mathbf{X}\) ; if, moreover, \(f\) tends to \(f_0\) in \(\mathcal{C}(\mathbf{X})\) , then \(\chi(s, \cdot)f\) tends to \(\chi(s_0, \cdot)f_0\) uniformly on every compact subset of \(\mathbf{X}\) , whence our assertion. Thus the representation \(\gamma_{\chi}\) of \(\mathbf{G}\) in \(\mathcal{C}(\mathbf{X})\) is continuous.

Let us show that the representation \(\pmb{\gamma}_{\chi}\) of G in \(\mathcal{K}(\mathrm{X})\) is continuous. Since \(\mathcal{K}(\mathrm{X})\) is the direct limit of Banach spaces, it is barreled (TVS, III, §4, No.~1, Cor. 3 of Prop. 3), thus it suffices to prove that \(\pmb{\gamma}_{\chi}\) is separately continuous (No.~1, Prop. 1). Now, let H be a compact subset of X and let \(s_0 \in \mathrm{G}\) . Let V be a compact neighborhood of \(s_0\) in G, and let L = VH, which is compact in X. For every \(f \in \mathcal{K}(\mathrm{X}, \mathrm{H})\) , the support of \(\pmb{\gamma}_{\chi}(s_0)f\) is contained in L, and

\[
\sup _ {x \in X} | (\boldsymbol {\gamma} _ {\chi} (s _ {0}) f) (x) | \leqslant \sup _ {x \in L} | \chi (s _ {0} ^ {- 1}, x) | \cdot \sup _ {x \in X} | f (x) |,
\]

therefore \(f \mapsto \boldsymbol{\gamma}_{\chi}(s_0)f\) is a continuous linear mapping of \(\mathcal{K}(\mathrm{X},\mathrm{H})\) into \(\mathcal{K}(\mathrm{X},\mathrm{L})\) ; it follows that \(f \mapsto \boldsymbol{\gamma}_{\chi}(s_0)f\) is a continuous linear mapping of \(\mathcal{K}(\mathrm{X})\) into itself (TVS, II, §4, No.~4, Prop. 5). On the other hand, the topology of \(\mathcal{K}(\mathrm{X},\mathrm{L})\) is induced by that of \(\mathcal{C}(\mathrm{X})\) . By what has already been proved, the mapping \(s \mapsto \boldsymbol{\gamma}_{\chi}(s)f\) of \(\mathrm{V}\) into \(\mathcal{K}(\mathrm{X},\mathrm{L})\) is continuous. This completes the proof that \(\boldsymbol{\gamma}_{\chi}\) is separately continuous.

PROPOSITION 5. --- Suppose that each function \(\chi(s,\cdot)\) is bounded. Then \(\gamma_{\chi}\) leaves \(\overline{\mathcal{K}(X)}\) stable, and the linear representation \(\gamma_{\chi}\) of \(G\) in \(\overline{\mathcal{K}(X)}\) is continuous.

It is clear that \(\pmb{\gamma}_{\chi}\) leaves \(\overline{\mathcal{K}(\mathrm{X})}\) stable and that each of the \(\pmb{\gamma}_{\chi}(s)\) is continuous in \(\overline{\mathcal{K}(\mathrm{X})}\) . On the other hand, for every \(f \in \mathcal{K}(\mathrm{X})\) , \(s \mapsto \pmb{\gamma}_{\chi}(s)f\) is a continuous mapping of \(\mathbf{G}\) into \(\mathcal{K}(\mathbf{X})\) and a fortiori into \(\overline{\mathcal{K}(\mathbf{X})}\) . Therefore the representation \(\pmb{\gamma}_{\chi}\) in \(\overline{\mathcal{K}(\mathbf{X})}\) is continuous (No.~1, Prop. 2).

